% =====================================================================
\section{Integrales de línea, superficie y volumen}
\label{sec:integrales}
% =====================================================================

\subsection{Integrales de línea: circulación y longitud}

\begin{proposicion}[Cálculo de una integral de línea]
\label{prop:linea}
Si $C$ está parametrizada por $\vb{r}(t)$, $t\in[t_{0},t_{1}]$, de $P$ a $Q$,
\begin{equation}
  \int_{C}\vb{A}\cdot\dd\vb{r}=\int_{t_{0}}^{t_{1}}\vb{A}\bigl(\vb{r}(t)\bigr)\cdot\dot{\vb{r}}(t)\,\dd t
  =\int_{P}^{Q}\bigl(A_{1}h_{1}\dd q^{1}+A_{2}h_{2}\dd q^{2}+A_{3}h_{3}\dd q^{3}\bigr),
  \label{eq:linea}
\end{equation}
la última forma en un sistema ortogonal, usando \eqref{eq:metrica}. La
integral \emph{depende del camino} salvo que $\vb{A}$ sea conservativo
(\cref{subsec:conservativos}).
\end{proposicion}

\begin{ejemplo}[Circulación sobre una circunferencia]
\label{ej:circulacion}
Calcular la circulación de $\vb{A}=r\sin\theta\,\er+2r\,\ef$ a lo largo de
la circunferencia $r=1$ del plano $XY$, recorrida en sentido antihorario.
\begin{enumerate}
  \item \emph{Desplazamiento}. En el plano $XY$ ($\theta=\pi/2$),
        $\dd\vb{r}=\dd r\,\er+r\,\dd\theta\,\et+r\sin\theta\,\dd\varphi\,\ef$ con
        $\dd r=\dd\theta=0$: $\dd\vb{r}=r\,\dd\varphi\,\ef$.
  \item \emph{Integrando}. $\vb{A}\cdot\dd\vb{r}=2r\cdot r\,\dd\varphi=2r^{2}\dd\varphi$
        (la componente radial no trabaja).
  \item \emph{Integral}. $\Gamma=\displaystyle\int_{0}^{2\pi}2r^{2}\dd\varphi=4\pi r^{2}\big|_{r=1}=4\pi$.
\end{enumerate}
\end{ejemplo}

\begin{proposicion}[Longitud de una curva]
\begin{equation}
  l_{PQ}=\int_{P}^{Q}\dd s=\int_{t_{0}}^{t_{1}}\sqrt{\dot x^{2}+\dot y^{2}+\dot z^{2}}\,\dd t
  =\int_{t_{0}}^{t_{1}}\sqrt{\textstyle\sum_{i}h_{i}^{2}(\dot q^{i})^{2}}\,\dd t .
  \label{eq:longitud}
\end{equation}
\end{proposicion}

\begin{ejemplo}[Hélice]
\label{ej:helice-longitud}
$x=A\cos\omega t$, $y=A\sin\omega t$, $z=vt$, $t\in[0,2\pi]$. Entonces
$\dot x^{2}+\dot y^{2}+\dot z^{2}=(A\omega)^{2}+v^{2}$, constante, y
$l=\int_{0}^{2\pi}\sqrt{(A\omega)^{2}+v^{2}}\,\dd t=2\pi\sqrt{(A\omega)^{2}+v^{2}}$.
\end{ejemplo}

\subsection{Integrales de superficie: flujo}

\begin{proposicion}[Cálculo de un flujo]
\label{prop:flujo}
Si $S$ está parametrizada por $\vb{r}(u,v)$, $(u,v)\in D$,
\begin{equation}
  \int_{S}\vb{A}\cdot\dd\vb{S}=\iint_{D}\vb{A}\cdot\Bigl(\pd{\vb{r}}{u}\times\pd{\vb{r}}{v}\Bigr)\dd u\,\dd v .
\end{equation}
Sobre una superficie coordenada $q^{3}=\text{const}$, por \eqref{eq:dS},
$\dd\vb{S}=h_{1}h_{2}\dd q^{1}\dd q^{2}\uv{e}_{3}$ y el flujo es la integral doble
\begin{equation}
  \Phi=\int_{a}^{b}\dd q^{1}\int_{c(q^{1})}^{d(q^{1})}\dd q^{2}\,\bigl[A_{3}h_{1}h_{2}\bigr](q^{1},q^{2},q^{3}).
\end{equation}
\end{proposicion}

\begin{teorema}[Fubini]
Si el integrando es continuo y el dominio es regular, las integrales
iteradas pueden calcularse en cualquier orden y coinciden con la integral
múltiple.
\end{teorema}

% ---------------------------------------------------------------------
\begin{figure}[ht]
\centering
\begin{tikzpicture}[scale=2.3]
  \fill[cian,opacity=0.08] (0,0) -- (1,0) arc[start angle=0,end angle=90,radius=1] -- cycle;
  \draw[eje] (0,0) -- (1.25,0) node[anchor=north,peq]{$x$};
  \draw[eje] (0,0) -- (0,1.2) node[anchor=east,peq]{$y$};
  \draw[texto,line width=0.9pt] (1,0) arc[start angle=0,end angle=90,radius=1];
  \pgfmathsetmacro{\yy}{0.55}\pgfmathsetmacro{\xx}{sqrt(1-\yy*\yy)}
  \fill[cian,opacity=0.45] (0,\yy) rectangle (\xx,\yy+0.07);
  \draw[cian,line width=0.5pt] (0,\yy) rectangle (\xx,\yy+0.07);
  \node[peq,anchor=east,inner sep=2pt] at (0,\yy+0.035) {$\dd y$};
  \draw[aux] (\xx,0) -- (\xx,\yy);
  \node[peq,anchor=north,inner sep=2pt] at (\xx,0) {$\sqrt{1-y^{2}}$};
  \node[peq,anchor=south west,inner sep=1pt] at ({cos(60)},{sin(60)}) {$x^{2}+y^{2}=1$};
  \node[etiq,anchor=south] at (0.55,1.25) {(a) Cartesianas};
\end{tikzpicture}\hspace{12mm}%
\begin{tikzpicture}[scale=2.3]
  \fill[violeta,opacity=0.08] (0,0) -- (1,0) arc[start angle=0,end angle=90,radius=1] -- cycle;
  \draw[eje] (0,0) -- (1.25,0) node[anchor=north,peq]{$x$};
  \draw[eje] (0,0) -- (0,1.2) node[anchor=east,peq]{$y$};
  \draw[texto,line width=0.9pt] (1,0) arc[start angle=0,end angle=90,radius=1];
  \def\ra{0.55}\def\rb{0.68}\def\fa{32}\def\fb{44}
  \fill[violeta,opacity=0.5] ({\ra*cos(\fa)},{\ra*sin(\fa)}) arc[start angle=\fa,end angle=\fb,radius=\ra]
     -- ({\rb*cos(\fb)},{\rb*sin(\fb)}) arc[start angle=\fb,end angle=\fa,radius=\rb] -- cycle;
  \draw[aux] (0,0) -- ({\rb*cos(\fa)},{\rb*sin(\fa)});
  \draw[aux] (0,0) -- ({\rb*cos(\fb)},{\rb*sin(\fb)});
  \node[peq,anchor=north west,inner sep=1pt] at ({0.62*cos(\fa)},{0.62*sin(\fa)}) {$\dd\rho$};
  \node[peq,anchor=south east,inner sep=1pt] at ({\rb*cos(\fb+2)},{\rb*sin(\fb+2)}) {$\rho\,\dd\varphi$};
  \node[etiq,anchor=south] at (0.55,1.25) {(b) Polares};
\end{tikzpicture}
\caption{Área del cuadrante de círculo unidad. (a) Franjas horizontales de
anchura $\sqrt{1-y^{2}}$ y espesor $\dd y$. (b) Elemento polar de lados
$\dd\rho$ y $\rho\,\dd\varphi$, de área $\rho\,\dd\rho\,\dd\varphi$.}
\label{fig:cuadrante}
\end{figure}
% ---------------------------------------------------------------------

\begin{ejemplo}[Área del cuadrante, \cref{fig:cuadrante}]
\label{ej:cuadrante}
\emph{Cartesianas}: con el cambio $y=\sin u$,
\[
  S=\int_{0}^{1}\dd y\int_{0}^{\sqrt{1-y^{2}}}\dd x=\int_{0}^{1}\sqrt{1-y^{2}}\,\dd y
  =\int_{0}^{\pi/2}\cos^{2}u\,\dd u=\frac{\pi}{4}.
\]
\emph{Polares} ($h_{\rho}h_{\varphi}=\rho$):
$S=\displaystyle\int_{0}^{1}\dd\rho\int_{0}^{\pi/2}\rho\,\dd\varphi=\frac{\pi}{2}\int_{0}^{1}\rho\,\dd\rho=\frac{\pi}{4}$.
Las coordenadas adaptadas a la frontera convierten los límites variables en constantes.
\end{ejemplo}

\subsection{Integrales de volumen}

\begin{teorema}[Cambio de variables]
\label{teo:cambio}
Para un sistema ortogonal, por \cref{prop:elementos},
\begin{equation}
  \int_{V}\Phi\,\dd V=\int\!\!\int\!\!\int\dd q^{1}\dd q^{2}\dd q^{3}\,\bigl[\Phi\,h_{1}h_{2}h_{3}\bigr](q^{1},q^{2},q^{3}),
\end{equation}
y, por Fubini, el orden de integración puede intercambiarse.
\end{teorema}

\begin{ejemplo}[Volumen de una semiesfera de radio $R$]
\emph{Esféricas} ($\dd V=r^{2}\sin\theta\,\dd r\,\dd\theta\,\dd\varphi$, límites constantes):
\[
  V=\int_{0}^{R}r^{2}\dd r\int_{0}^{\pi/2}\sin\theta\,\dd\theta\int_{0}^{2\pi}\dd\varphi
  =\Bigl[\tfrac{r^{3}}{3}\Bigr]_{0}^{R}\bigl[-\cos\theta\bigr]_{0}^{\pi/2}\bigl[\varphi\bigr]_{0}^{2\pi}
  =\frac{2\pi}{3}R^{3}.
\]
\emph{Cartesianas}: por simetría basta $\tfrac14$ de la semiesfera. Con
$a\defeq\sqrt{R^{2}-x^{2}}$,
\[
  V=4\int_{0}^{R}\dd x\int_{0}^{a}\dd y\int_{0}^{\sqrt{a^{2}-y^{2}}}\dd z
  =4\int_{0}^{R}\dd x\underbrace{\int_{0}^{a}\sqrt{a^{2}-y^{2}}\,\dd y}_{\pi a^{2}/4\ (\cref{ej:cuadrante})}
  =\pi\int_{0}^{R}(R^{2}-x^{2})\,\dd x=\frac{2\pi}{3}R^{3}.
\]
\end{ejemplo}

\subsection{Teoremas de Gauss y de Stokes}

\begin{teorema}[Gauss--Ostrogradski, de la divergencia]
\label{teo:gauss}
Sea $V\subset\R^{3}$ una región acotada cuya frontera $\partial V$ es una
superficie cerrada regular a trozos, orientada con la normal saliente, y
$\vb{A}$ un campo de clase $\C^{1}$ en un abierto que contiene a $V$. Entonces
\begin{equation}
  \resultado{\oint_{\partial V}\vb{A}\cdot\dd\vb{S}=\int_{V}\div\vb{A}\,\dd V}
  \label{eq:gauss}
\end{equation}
\end{teorema}

\begin{teorema}[Stokes]
\label{teo:stokes}
Sea $S$ una superficie orientable regular a trozos cuya frontera
$C=\partial S$ es una curva cerrada, orientada con $S$ por la regla de la
mano derecha, y $\vb{A}$ de clase $\C^{1}$ en un abierto que contiene a $S$. Entonces
\begin{equation}
  \resultado{\oint_{\partial S}\vb{A}\cdot\dd\vb{r}=\int_{S}(\rot\vb{A})\cdot\dd\vb{S}}
  \label{eq:stokes}
\end{equation}
\end{teorema}

% ---------------------------------------------------------------------
\begin{figure}[ht]
\centering
\begin{tikzpicture}[scale=1.0,baseline=(tit.base)]
  % Stokes: casquete sobre una curva cerrada
  \shade[inner color=violeta!22!caja,outer color=caja]
     (-2,0) arc[start angle=180,end angle=360,x radius=2,y radius=0.6]
     .. controls (2,1.6) and (1,2.4) .. (0,2.4) .. controls (-1,2.4) and (-2,1.6) .. (-2,0);
  \draw[texto!70,line width=0.5pt] (2,0) .. controls (2,1.6) and (1,2.4) .. (0,2.4)
     .. controls (-1,2.4) and (-2,1.6) .. (-2,0);
  \draw[cian,line width=0.6pt,dash pattern=on 2pt off 2pt] (2,0) arc[start angle=0,end angle=180,x radius=2,y radius=0.6];
  \draw[cian,line width=1.1pt,postaction={decorate},
    decoration={markings,mark=at position 0.5 with {\arrow{Stealth[length=7pt]}}}]
    (-2,0) arc[start angle=180,end angle=360,x radius=2,y radius=0.6];
  \node[cian,font=\small,anchor=north] at (0,-0.65) {$C=\partial S$};
  \draw[vec,rosa] (0,2.4) -- (0,3.25) node[anchor=south,etiq,inner sep=1pt]{$\un$};
  \draw[vec,rosa] (1.45,1.55) -- (2.1,2.1);
  \draw[vec,rosa] (-1.45,1.55) -- (-2.1,2.1);
  \node[violeta,font=\small] at (0,1.25) {$S$};
  \node[etiq,anchor=south] (tit) at (0,3.75) {(a) Stokes};
\end{tikzpicture}\hspace{14mm}%
\begin{tikzpicture}[scale=1.0,baseline=(tit.base)]
  % Gauss: elipsoide de revolución con normales salientes exactas
  \def\a{2.0}\def\b{1.45}\def\cy{1.25}
  \shade[ball color=cian!40!caja,opacity=0.55] (0,\cy) ellipse[x radius=\a,y radius=\b];
  \fill[fondo,opacity=0.25] (0,\cy) ellipse[x radius=\a,y radius=\b];
  % ecuador: mitad trasera a trazos, delantera continua
  \draw[cian!70,line width=0.5pt,dash pattern=on 2pt off 2pt]
     (\a,\cy) arc[start angle=0,end angle=180,x radius=\a,y radius=0.42];
  \draw[cian!70,line width=0.6pt]
     (-\a,\cy) arc[start angle=180,end angle=360,x radius=\a,y radius=0.42];
  % silueta
  \draw[texto!80,line width=0.7pt] (0,\cy) ellipse[x radius=\a,y radius=\b];
  % elemento dS resaltado sobre la silueta
  \draw[rosa,line width=2.2pt] plot[domain=50:70,samples=12,smooth]
     ({\a*cos(\x)},{\cy+\b*sin(\x)});
  % normales salientes: n ∝ (b cos t, a sin t)
  \foreach \t/\len in {60/0.95,140/0.7,200/0.7,265/0.7,320/0.7,10/0.7}{
    \pgfmathsetmacro{\nx}{\b*cos(\t)}\pgfmathsetmacro{\ny}{\a*sin(\t)}
    \pgfmathsetmacro{\nn}{sqrt(\nx*\nx+\ny*\ny)}
    \draw[vec,rosa] ({\a*cos(\t)},{\cy+\b*sin(\t)}) -- ++({\len*\nx/\nn},{\len*\ny/\nn}) coordinate (N\t);}
  \node[etiq,anchor=south west,inner sep=1pt] at (N60) {$\dd\vb{S}=\un\,\dd S$};
  \node[cian,font=\small] at (0,\cy) {$V$};
  \node[texto!80,font=\small,anchor=south east,inner sep=2pt] at ({\a*cos(108)},{\cy+\b*sin(108)}) {$\partial V$};
  \node[etiq,anchor=south] (tit) at (0,3.75) {(b) Gauss};
\end{tikzpicture}
\caption{(a) Teorema de Stokes: la orientación de $C$ y la normal $\un$ de $S$
están ligadas por la regla de la mano derecha; cualquier superficie con el
mismo borde da el mismo flujo de $\rot\vb{A}$. (b) Teorema de Gauss: la
frontera $\partial V$ se orienta con la normal \emph{saliente}.}
\label{fig:gauss-stokes}
\end{figure}
% ---------------------------------------------------------------------

\begin{observacion}[Lectura local]
Aplicados a un volumen o a una superficie infinitesimales, \eqref{eq:gauss}
y \eqref{eq:stokes} \emph{son} las definiciones \eqref{eq:div-def} y
\eqref{eq:rot-def}: la divergencia es flujo por unidad de volumen y el
rotacional, circulación por unidad de área.
\end{observacion}

\begin{observacion}[Las identidades nulas, vistas por Stokes y Gauss]
Para una superficie \emph{cerrada} $\partial V$ no hay borde, luego
$\oint_{\partial V}(\rot\vb{A})\cdot\dd\vb{S}=0$ y, por Gauss,
$\int_{V}\div(\rot\vb{A})\,\dd V=0$ para todo $V$: de nuevo
$\div(\rot\vb{A})=0$ (\cref{teo:nulas}).
\end{observacion}

\subsection{Campos conservativos}
\label{subsec:conservativos}

\begin{definicion}[Campo conservativo]
$\vb{A}$ es \emph{conservativo} en $D$ si su circulación entre dos puntos
no depende del camino:
$\int_{P,C_{1}}^{Q}\vb{A}\cdot\dd\vb{r}=\int_{P,C_{2}}^{Q}\vb{A}\cdot\dd\vb{r}$
para cualesquiera $C_{1},C_{2}$; equivalentemente, $\oint_{C}\vb{A}\cdot\dd\vb{r}=0$
para toda curva cerrada $C=C_{1}\cup(-C_{2})$.
\end{definicion}

\begin{teorema}[Caracterización de los campos conservativos]
\label{teo:conservativo}
Sea $\vb{A}\in\C^{1}(D)$ con $D\subseteq\R^{3}$ abierto y conexo. Son
equivalentes:
\begin{casos}
  \item $\vb{A}$ es conservativo;
  \item existe $\Phi\in\C^{2}(D)$ (potencial) con $\vb{A}=\grad\Phi$.
\end{casos}
En tal caso $\rot\vb{A}=\vb{0}$. Si además $D$ es \emph{simplemente conexo},
también $\rot\vb{A}=\vb{0}\ \Longrightarrow\ \vb{A}$ conservativo.
\end{teorema}

\begin{proof}
(ii)$\Rightarrow$(i): $\int_{P}^{Q}\grad\Phi\cdot\dd\vb{r}=\int_{P}^{Q}\dd\Phi=\Phi(Q)-\Phi(P)$,
independiente del camino. (i)$\Rightarrow$(ii): fijado $P_{0}$, se define
$\Phi(\vb{r})=\int_{P_{0}}^{\vb{r}}\vb{A}\cdot\dd\vb{r}'$ (bien definida por (i)); entonces
$\dd\Phi=\vb{A}\cdot\dd\vb{r}$ y, por \eqref{eq:grad-def}, $\vb{A}=\grad\Phi$.
$\rot\vb{A}=\rot\grad\Phi=\vb{0}$ por \cref{teo:nulas}. Recíprocamente, si
$D$ es simplemente conexo toda curva cerrada $C\subset D$ es borde de una
superficie $S\subset D$ y, por Stokes,
$\oint_{C}\vb{A}\cdot\dd\vb{r}=\int_{S}(\rot\vb{A})\cdot\dd\vb{S}=0$.
\end{proof}

\begin{observacion}[La hipótesis topológica importa]
$\vb{A}=\dfrac{-y\,\ii+x\,\jj}{x^{2}+y^{2}}=\dfrac{\ef}{\rho}$ cumple
$\rot\vb{A}=\vb{0}$ en $D=\R^{3}\setminus\{\text{eje }z\}$ (por \eqref{eq:rot-cil}:
$\rho A_{\varphi}=1$ es constante), pero su circulación sobre la circunferencia
$\rho=1$ es $\int_{0}^{2\pi}\tfrac{1}{\rho}\,\rho\,\dd\varphi=2\pi\neq0$: $D$ no es simplemente conexo.
\end{observacion}

\begin{clave}{Método: hallar el potencial de $\vb{A}$}
\begin{enumerate}[leftmargin=1.4em]
  \item \emph{Comprobar} $\rot\vb{A}=\vb{0}$ en un dominio simplemente conexo
        (en $\R^{3}$ basta).
  \item \emph{Integración sucesiva}: $\Phi=\int A_{x}\,\dd x+g(y,z)$; imponer
        $\partial_{y}\Phi=A_{y}$ para hallar $g$ salvo $h(z)$; imponer
        $\partial_{z}\Phi=A_{z}$ para hallar $h$ salvo constante.
  \item \emph{Alternativa}: integrar por la poligonal
        $(0,0,0)\to(x,0,0)\to(x,y,0)\to(x,y,z)$:
        \begin{equation}
          \Phi(x,y,z)=\int_{0}^{x}\!A_{x}(s,0,0)\dd s+\int_{0}^{y}\!A_{y}(x,s,0)\dd s+\int_{0}^{z}\!A_{z}(x,y,s)\dd s .
          \label{eq:poligonal}
        \end{equation}
  \item \emph{Verificar} $\grad\Phi=\vb{A}$.
\end{enumerate}
\end{clave}
